% Part: propositional-logic
% Chapter: syntax-and-semantics
% Section: formulas

\documentclass[../../../include/open-logic-section]{subfiles}

\begin{document}

\olfileid[hi]{pl}{syn}{fml}

\olsection{प्रतिज्ञप्ति \usetoken{S}{formula}}

प्रतिज्ञप्ति तर्कशास्त्र में !!^{formula}s का निर्माण \emph{!!{propositional
    variable}s}\iftag{prvFalse}{\iftag{prvTrue}{,}{ तथा} प्रतिज्ञप्ति
  नियतांक~$\lfalse$}{}\iftag{prvTrue}{\iftag{prvFalse}{ तथा}{} प्रतिज्ञप्ति
  नियतांक~$\ltrue$}{} और \emph{तार्किक संयोजकों} से होता है।

\begin{enumerate}
\item !!^a{denumerable} समुच्चय~$\PVar$, जिसमें !!{propositional variable}s
  $\Obj p_0$, $\Obj p_1$, \dots हैं।
\tagitem{prvFalse}{!!{falsity} के लिए प्रतिज्ञप्ति नियतांक~$\lfalse$।}{}
\tagitem{prvTrue}{!!{truth} के लिए प्रतिज्ञप्ति नियतांक~$\ltrue$।}{}
\item तार्किक संयोजक:
  \startycommalist
  \iftag{prvNot}{\ycomma $\lnot$ (निषेध)}{}%
  \iftag{prvAnd}{\ycomma $\land$ (संयोजन)}{}%
  \iftag{prvOr}{\ycomma $\lor$ (वियोजन)}{}%
  \iftag{prvIf}{\ycomma $\lif$ (!!{conditional})}{}%
  \iftag{prvIff}{\ycomma $\liff$ (!!{biconditional})}{}%
\item विराम-चिह्न: (, ), और अल्पविराम।
\end{enumerate}

प्रतिज्ञप्ति तर्कशास्त्र की इस भाषा को हम $\Lang L_0$ से निरूपित करते हैं।

\iftag{defNot,defOr,defAnd,defIf,defIff,defTrue,defFalse,defEx,defAll}{%
ऊपर प्रस्तुत आदिम संयोजकों के अतिरिक्त हम निम्नलिखित \emph{परिभाषित}
प्रतीकों का भी उपयोग करते हैं:
\startycommalist
  \iftag{defNot}{\ycomma $\lnot$ (निषेध)}{}%
  \iftag{defAnd}{\ycomma $\land$ (संयोजन)}{}%
  \iftag{defOr}{\ycomma $\lor$ (वियोजन)}{}%
  \iftag{defIf}{\ycomma $\lif$ (!!{conditional})}{}%
  \iftag{defIff}{\ycomma $\liff$ (!!{biconditional})}{}%
  \iftag{defFalse}{\ycomma $\lfalse$ (!!{falsity})}{}%
  \iftag{defTrue}{\ycomma $\ltrue$ (!!{truth})}{}.}{}

\begin{tagblock}{defNot,defOr,defAnd,defIf,defIff,defTrue,defFalse,defEx,defAll}
\begin{explain}
परिभाषित प्रतीक औपचारिक रूप से भाषा का भाग नहीं होता, बल्कि उसे अनौपचारिक
संक्षेप के रूप में प्रस्तुत किया जाता है: उससे हम उन सूत्रों को संक्षिप्त लिख
सकते हैं जो केवल आदिम प्रतीकों के प्रयोग पर बहुत लंबे हो जाते। यह स्पष्टतः
लाभदायक है। किंतु बड़ा लाभ यह है कि प्रमाण छोटे हो जाते हैं। यदि कोई प्रतीक
आदिम हो, तो प्रमाणों में उसे अलग से संभालना पड़ता है। अतः आदिम प्रतीक जितने
अधिक होंगे, हमारे प्रमाण उतने ही लंबे होंगे।
\end{explain}
\end{tagblock}

% Alternate symbols

\begin{intro}
हो सकता है कि आप ऊपर प्रयुक्त पदावली और प्रतीकों से भिन्न रूपों से परिचित
हों। तर्कशास्त्र की पुस्तकों (और शिक्षकों) में ``निषेध'' के लिए प्रायः
$\sim$, $\neg$, और~!\ तथा ``संयोजन'' के लिए $\wedge$, $\cdot$, और $\&$
प्रयुक्त होते हैं। ``आपादन'' या ``निहितार्थ'' के लिए सामान्यतः
$\rightarrow$, $\Rightarrow$, और $\supset$ प्रतीक प्रयुक्त होते हैं।
\iftag{prvIff,defIff}{``द्विशर्त,'' ``द्वि-आपादन,'' या ``(शाब्दिक)
  तुल्यता'' के प्रतीक $\leftrightarrow$,
  $\Leftrightarrow$, और $\equiv$ हैं।}{}
\iftag{prvFalse,defFalse}{$\lfalse$ प्रतीक को अलग-अलग प्रसंगों में
  ``असत्यता,'' ``असत्य,'' ``असंगति,'' या ``तल'' कहा जाता है।}{}
\iftag{prvTrue,defTrue}{$\ltrue$ प्रतीक को अलग-अलग प्रसंगों में
  ``सत्यता,'' ``सत्य,'' या ``शीर्ष'' कहा जाता है।}{}
\end{intro}

\begin{defn}[सूत्र]
\ollabel{defn:formulas}
प्रतिज्ञप्ति तर्कशास्त्र के \emph{!!{formula}s} का समुच्चय~$\Frm[L_0]$
निम्न प्रकार आगमनात्मक रूप से परिभाषित है:
\begin{enumerate}
\tagitem{prvFalse}{$\lfalse$ एक परमाणु !!{formula} है।}{}

\tagitem{prvTrue}{$\ltrue$ एक परमाणु !!{formula} है।}{}

\item प्रत्येक !!{propositional variable}~$\Obj p_i$ एक परमाणु
  !!{formula} है।

\tagitem{prvNot}{यदि $!A$ !!a{formula} है, तो $\lnot !A$ भी
  !!a{formula} है।}{}

\tagitem{prvAnd}{यदि $!A$ और $!B$ !!{formula}s में से हों, तो $(!A \land
  !B)$ भी !!a{formula} है।}{}

\tagitem{prvOr}{यदि $!A$ और $!B$ !!{formula}s में से हों, तो $(!A \lor !B)$
  भी !!a{formula} है।}{}

\tagitem{prvIf}{यदि $!A$ और $!B$ !!{formula}s में से हों, तो $(!A \lif !B)$
  भी !!a{formula} है।}{}

\tagitem{prvIff}{यदि $!A$ और $!B$ !!{formula}s में से हों, तो $(!A \liff !B)$
  भी !!a{formula} है।}{}

\tagitem{limitClause}{इसके अतिरिक्त कुछ भी !!a{formula} नहीं है।}{}
\end{enumerate}
\end{defn}

\begin{explain}
!!{formula}s की यह परिभाषा एक \emph{आगमनात्मक परिभाषा} है। मूलतः हम
!!{formula}s का समुच्चय अनंत चरणों में बनाते हैं। आरंभिक चरण में हम सभी
परमाणु सूत्रों को सूत्र घोषित करते हैं; यह परिभाषा के आरंभिक मामलों—अर्थात्
इनके लिए दिए मामलों—के अनुरूप है:
\iftag{prvTrue}{$\ltrue$, }{}%
\iftag{prvFalse}{$\lfalse$, }{}%
$\Obj p_i$. अतः ``परमाण्विक !!{formula}'' का अर्थ इस रूप का कोई भी
!!{formula} है।

परिभाषा के अन्य मामले नए !!{formula}s के निर्माण के नियम देते हैं, जिनमें
पहले से निर्मित !!{formula}s का उपयोग होता है। दूसरे चरण में हम इन नियमों
से !!{formula}s का निर्माण परमाणु !!{formula}s से कर सकते हैं। तीसरे चरण में हम परमाणु
सूत्रों और दूसरे चरण में प्राप्त सूत्रों से नए सूत्र बनाते हैं, और यह क्रम
आगे चलता है। !!{formula} वही और केवल वही है जो अंततः ऐसे किसी चरण में
बनाया जाता है।
\end{explain}

सूत्र $(!B \ast !C)$ को लिखते समय, जो $!B$, $!C$ से दो-स्थानी
संयोजक~$\ast$ का उपयोग करके बना है, हम प्रायः सबसे बाहरी कोष्ठक-युग्म
छोड़कर केवल~$!B \ast !C$ लिखेंगे।

\begin{tagblock}{defNot,defOr,defAnd,defIf,defIff,defTrue,defFalse,defEx,defAll}
\begin{defn}
परिभाषित संक्रियकों से बने सूत्रों को इस प्रकार समझना है:

\begin{tagenumerate}{defTrue,defFalse,defNot,defOr,defAnd,defIf,defIff,defEx,defAll}

\tagitem{defTrue}{$\ltrue$ इसका संक्षेप है:
  \iftag{prvFalse}{$\lnot\lfalse$}{$(!A \lor \lnot !A)$, जहाँ कोई
    नियत परमाणु !!{formula}~$!A$ लिया गया है।}.}{}

\tagitem{defFalse}{$\lfalse$ इसका संक्षेप है:
  \iftag{prvTrue}{$\lnot\ltrue$}{$(!A \land \lnot !A)$, जहाँ कोई
    नियत परमाणु !!{formula}~$!A$ लिया गया है।}.}{}

\tagitem{defNot}{$\lnot !A$, $!A \lif \lfalse$ का संक्षेप है।}{}

\tagitem{defOr}{$!A \lor !B$ इसका संक्षेप है:
  \iftag{prvAnd}{$\lnot(\lnot !A \land \lnot !B)$}{$\lnot !A \lif
    !B$।}.}{}

\tagitem{defAnd}{$!A \land !B$ इसका संक्षेप है:
  \iftag{prvOr}{$\lnot(\lnot !A \lor \lnot !B)$}{$\lnot (!A \lif
    \lnot !B)$।}.}{}

\tagitem{defIf}{$!A \lif !B$ इसका संक्षेप है:
  \iftag{prvOr}{$\lnot !A \lor !B)$}{$\lnot (!A \land \lnot !B)$}.}{}

\tagitem{defIff}{$!A \liff !B$, $(!A \lif !B) \land (!B
  \lif !A)$ का संक्षेप है।}{}
\end{tagenumerate}
\end{defn}
\end{tagblock}

\begin{defn}[वाक्यविन्यासगत समानता]
प्रतीक $\ident$ बताता है कि दो प्रतीक-शृंखलाएँ विन्यास में समान हैं;
अर्थात् $!A \ident !B$ तभी और केवल तभी, जब $!A$ और $!B$ समान लंबाई की
प्रतीक-शृंखलाएँ हों और प्रत्येक स्थान पर उनमें एक ही प्रतीक हो।
\end{defn}

प्रतीक $\ident$ के दोनों ओर संयोजन से प्राप्त शृंखलाएँ रखी जा सकती हैं।
उदाहरणतः, $!A \ident (!B \lor !C)$ का अर्थ है: प्रतीक-शृंखला~$!A$ वही
शृंखला है जो इसी क्रम में आरंभिक कोष्ठक, शृंखला $!B$, प्रतीक $\lor$,
शृंखला~$!C$, और अंतिम कोष्ठक को जोड़ने से प्राप्त होती है। यदि ऐसा है, तो
हम जानते हैं कि $!A$ का पहला प्रतीक आरंभिक कोष्ठक है, $!A$ में $!B$ एक
उपशृंखला के रूप में (दूसरे प्रतीक से आरंभ होकर) आता है, उसके बाद $\lor$
आता है, आदि।

\end{document}
